\documentclass[12pt]{letter}
\usepackage[margin=1in]{geometry}
\usepackage{hyperref}

\signature{Xinze Li}
\address{School of Mechanical Engineering \\
[University Name] \\
[City], [Country] \\
E-mail: [author@email.com]}

\begin{document}

\begin{letter}{Editor-in-Chief \\
IEEE Transactions on Instrumentation and Measurement}

\opening{Dear Editor-in-Chief,}

I am pleased to submit our manuscript entitled \textbf{``Perspective-Corrected
Monocular Vision for Non-Contact Motor Shaft Vibration Measurement''} for
consideration for publication in the IEEE Transactions on Instrumentation
and Measurement.

\section*{Relevance to IEEE TIM}

This manuscript directly aligns with IEEE TIM's scope in instrumentation and
measurement methodologies. We present a novel non-contact vibration measurement
system that uses a single monocular camera and visual object tracking to
measure motor shaft vibration amplitude. The work addresses a practical
measurement challenge---camera misalignment in field deployments---through a
rigorous geometric derivation and systematic experimental validation.

\section*{Novelty and Contribution}

The key innovation is a \textbf{perspective-corrected isotropic scale factor}
$s = D a / (2b^2)$, derived from the pinhole camera model, which correctly
converts pixel displacements to physical displacements when the camera views
a circular shaft cross-section obliquely. We prove that the scale factor is
isotropic (identical in all directions) despite the elliptical appearance of the
projected circle, and show that the commonly used naive scale $D/(2a)$
systematically underestimates the true scale by $\cos^2\theta$.

The method requires only the shaft diameter as prior knowledge---no depth
sensors, laser rangefinders, or target markers are needed. Validation on 640
synthetic images (8 tilt angles $\times$ 5 vibration phases $\times$ 4 noise
levels $\times$ 4 blur levels) and real motor shaft videos at 10--100 RPM
demonstrates the method's accuracy and robustness.

\section*{Key Findings}

\begin{itemize}
\item Tilt angle estimation accuracy: mean error $< 1^{\circ}$ for
$\theta \geq 3^{\circ}$
\item Scale formula error: $< 0.1\%$ (negligible compared to detection errors)
\item Dominant error source: ellipse center detection precision ($\sim 0.1$~px),
not the scale formula
\item Real-world improvement: scale stability improved by 2.9$\times$, drift
reduced by up to 3$\times$
\end{itemize}

\section*{Statements}

This manuscript has not been published previously and is not under consideration
for publication elsewhere. All authors have approved the manuscript and agree
with its submission to IEEE TIM. There are no conflicts of interest to declare.

The undergraduate thesis on which this work is based has been submitted to
[University Name] as a degree requirement. The thesis and this manuscript
contain distinct content, with this manuscript focusing on the perspective
correction methodology and including new analysis not present in the thesis.

Thank you for considering our manuscript. We look forward to your response.

\closing{Sincerely,}

\end{letter}
\end{document}
